\chapter{Chọn Nghề Và Áp Lực Định Hướng}
\markboth{Chọn Nghề Và Áp Lực Định Hướng}{Career Choice and Direction Pressure}

\section{Nhiều Học Sinh Chọn Nghề Để Tránh Bị Chê}
\begin{truyen}{Nhiều Học Sinh Chọn Nghề Để Tránh Bị Chê}{Many Students Choose Careers to Avoid Criticism}
\chuhoa{K}{hi được hỏi muốn học gì, không ít học sinh trả lời bằng thứ người khác sẽ nghĩ, không phải thứ mình thực sự muốn.}
\textit{(When asked what they want to study, many students answer with what others will approve of, not what they truly want.)}

Các em sợ học ngành bị coi là thấp, sợ họ hàng hỏi thu nhập, sợ bạn bè khoe trường danh giá hơn.
\textit{(They fear entering a field considered low status, fear relatives asking about income, fear classmates boasting about more prestigious schools.)}

Thế là các em chọn một con đường đủ an toàn để khỏi bị chê, rồi mang nỗi chán đó đi suốt nhiều năm.
\textit{(So they choose a path safe enough to avoid criticism, then carry that deadness for years.)}

Lựa chọn nghề nghiệp mà chỉ dựa trên ánh mắt thiên hạ rất dễ biến cuộc đời thành một màn trình diễn dài.
\textit{(A career choice based only on public opinion easily turns life into one long performance.)}

\end{truyen}

\begin{baihoc}
Định hướng nghề nghiệp phải bắt đầu từ năng lực, khí chất và giá trị sống, không chỉ từ danh tiếng bề ngoài.
\textit{(Career guidance must begin with ability, temperament, and values, not only external prestige.)}
\end{baihoc}

\begin{ghinhoanh}
\textbf{Short English to remember:}

Career guidance must begin with ability, temperament, and values, not only external prestige.

\end{ghinhoanh}

\begin{tuhocnhanh}
\textbf{Cau hoi tu hoc / Self-study questions}

1. Van de chinh la gi? \textit{What is the main problem?}

2. Cach xu ly tot hon la gi? \textit{What is a better action?}

3. Mot cau tieng Anh em muon nho la cau nao? \textit{Which English sentence do you want to remember?}
\end{tuhocnhanh}
\ngancach

\section{Đam Mê Không Đủ, Nhưng Không Có Nó Cũng Rất Mệt}
\begin{truyen}{Đam Mê Không Đủ, Nhưng Không Có Nó Cũng Rất Mệt}{Passion Is Not Enough, But Having None Is Exhausting}
\chuhoa{N}{gười lớn hay nói: cứ có đam mê là được.}
\textit{(Adults often say: passion is enough.)}

Câu đó đẹp nhưng thiếu.
\textit{(It sounds nice, but it is incomplete.)}

Đam mê không thay thế được năng lực, kỷ luật, và nhu cầu thị trường.
\textit{(Passion does not replace ability, discipline, and market reality.)}

Nhưng ngược lại, học một ngành mình hoàn toàn không thấy ý nghĩa cũng là một kiểu mòn mỏi rất thật.
\textit{(But the opposite is also true: studying a field you find no meaning in creates a very real kind of erosion.)}

Học sinh cần được dạy một phiên bản trưởng thành hơn: hãy tìm giao điểm giữa điều mình có thể làm tốt, điều mình thấy đáng làm, và điều xã hội cần.
\textit{(Students need a more mature version of the advice: find the overlap between what you can do well, what you find worth doing, and what society needs.)}

\end{truyen}

\begin{baihoc}
Định hướng tốt không thần thánh hóa đam mê, nhưng cũng không chà nát nó. Nó giúp học sinh nhìn nghề bằng cả trái tim lẫn đầu óc.
\textit{(Good guidance does not worship passion, but it does not crush it either. It helps students view work with both heart and mind.)}
\end{baihoc}

\begin{ghinhoanh}
\textbf{Short English to remember:}

Good guidance does not worship passion, but it does not crush it either.

It helps students view work with both heart and mind.

\end{ghinhoanh}

\begin{tuhocnhanh}
\textbf{Cau hoi tu hoc / Self-study questions}

1. Van de chinh la gi? \textit{What is the main problem?}

2. Cach xu ly tot hon la gi? \textit{What is a better action?}

3. Mot cau tieng Anh em muon nho la cau nao? \textit{Which English sentence do you want to remember?}
\end{tuhocnhanh}
\ngancach

\section{Không Ai Mười Bảy Tuổi Mà Hiểu Hết Mình}
\begin{truyen}{Không Ai Mười Bảy Tuổi Mà Hiểu Hết Mình}{No One Fully Understands Themselves at Seventeen}
\chuhoa{N}{hiều học sinh lo vì mình chưa biết chắc muốn gì.}
\textit{(Many students panic because they are not sure what they want.)}

Các em tưởng chỉ có mình mơ hồ.
\textit{(They think only they are uncertain.)}

Sự thật là phần lớn người lớn cũng từng mù mờ ở tuổi đó.
\textit{(The truth is that most adults were also unclear at that age.)}

Nhiều người ba mươi vẫn đang chỉnh lại hướng đi.
\textit{(Many people are still adjusting course at thirty.)}

Vấn đề không phải em chưa biết hết.
\textit{(The problem is not that you do not know everything.)}

Vấn đề là em có đang chịu khám phá nghiêm túc hay không.
\textit{(The problem is whether you are willing to explore seriously.)}

Thử việc nhỏ, đọc sâu về nghề, nói chuyện với người trong ngành, quan sát chính mình khi làm một việc cụ thể.
\textit{(Try small tasks, read deeply about fields, talk to people doing the work, observe yourself doing something concrete.)}

Đó là cách biết dần, không phải ngồi chờ một tiếng gọi huyền bí.
\textit{(That is how clarity grows, not by waiting for a mystical calling.)}

\end{truyen}

\begin{baihoc}
Hãy giảm áp lực phải biết ngay tất cả. Điều học sinh cần là tiến gần hơn tới câu trả lời bằng hành động thật.
\textit{(Reduce the pressure to know everything immediately. What students need is to move closer to answers through real action.)}
\end{baihoc}

\begin{ghinhoanh}
\textbf{Short English to remember:}

Reduce the pressure to know everything immediately.

What students need is to move closer to answers through real action.

\end{ghinhoanh}

\begin{tuhocnhanh}
\textbf{Cau hoi tu hoc / Self-study questions}

1. Van de chinh la gi? \textit{What is the main problem?}

2. Cach xu ly tot hon la gi? \textit{What is a better action?}

3. Mot cau tieng Anh em muon nho la cau nao? \textit{Which English sentence do you want to remember?}
\end{tuhocnhanh}
\ngancach

\section{Đừng Giao Cả Cuộc Đời Cho Một Trào Lưu}
\begin{truyen}{Đừng Giao Cả Cuộc Đời Cho Một Trào Lưu}{Do Not Hand Your Life to a Trend}
\chuhoa{C}{ứ vài năm lại có một nghề được tung hô như lối tắt thành công.}
\textit{(Every few years, a profession is marketed as the shortcut to success.)}

Học sinh dễ bị cuốn vào những câu như nghề hot, lương cao, tương lai sáng mà quên hỏi: mình có hợp không, có chịu được không, có muốn sống kiểu đó lâu dài không.
\textit{(Students are easily pulled in by phrases like hot field, high salary, bright future, and forget to ask: does it suit me, can I endure it, do I want that lifestyle long-term?)}

Một nghề có thể đang được thị trường săn đón, nhưng nếu nó làm em kiệt quệ mỗi ngày, cái giá sẽ rất đắt.
\textit{(A job may be highly demanded by the market, but if it drains you every day, the price will be high.)}

Chọn nghề theo trào lưu có thể giúp em đỡ bị lạc trong ngắn hạn, nhưng dễ làm em lạc sâu hơn về sau.
\textit{(Choosing by trend may reduce uncertainty in the short term, but can make you more lost later.)}

\end{truyen}

\begin{baihoc}
Học sinh cần được dạy cách nhìn xa hơn danh sách nghề hot. Nghề đáng theo là nghề em có thể lớn lên cùng nó, không chỉ chạy theo nó.
\textit{(Students need to be taught to look beyond lists of trendy jobs. A worthy profession is one you can grow with, not merely chase.)}
\end{baihoc}

\begin{ghinhoanh}
\textbf{Short English to remember:}

Students need to be taught to look beyond lists of trendy jobs.

A worthy profession is one you can grow with, not merely chase.

\end{ghinhoanh}

\begin{tuhocnhanh}
\textbf{Cau hoi tu hoc / Self-study questions}

1. Van de chinh la gi? \textit{What is the main problem?}

2. Cach xu ly tot hon la gi? \textit{What is a better action?}

3. Mot cau tieng Anh em muon nho la cau nao? \textit{Which English sentence do you want to remember?}
\end{tuhocnhanh}
\ngancach

\section{Con Nhà Thợ Rèn Không Nhất Thiết Phải Làm Thợ Rèn}
\begin{truyen}{Con Nhà Thợ Rèn Không Nhất Thiết Phải Làm Thợ Rèn}{The Blacksmith’s Child Need Not Become a Blacksmith}
\chuhoa{N}{gày xưa nghề nghiệp thường đi theo dòng họ.}
\textit{(Long ago professions often followed bloodlines.)}

Con nhà thợ rèn thành thợ rèn, con nhà làm mộc thành thợ mộc.
\textit{(A blacksmith’s child became a blacksmith, a carpenter’s child became a carpenter.)}

Cách đó có mặt tiện: trẻ được học nghề sớm.
\textit{(That system had an advantage: children learned skills early.)}

Nhưng nó cũng khóa luôn khả năng một đứa trẻ có thể hợp việc khác hơn.
\textit{(But it also locked away the possibility that a child might fit another calling better.)}

Một cậu bé mê chữ hơn mê búa bị coi là lệch.
\textit{(A boy who loved words more than hammers was considered strange.)}

Sau này chính cậu trở thành người ghi chép sổ sách giúp cả làng buôn bán minh bạch hơn.
\textit{(Later he became the record-keeper who helped the village trade more honestly.)}

Từ xưa đến nay, xã hội vẫn thường thích đường quen.
\textit{(Across time, society prefers familiar paths.)}

Nhưng phát triển thật lại cần chừa chỗ cho đường hợp.
\textit{(Real development, however, requires room for fitting paths.)}

\end{truyen}

\begin{baihoc}
Định hướng nghề nghiệp không nên chỉ dựa vào truyền thống gia đình. Hợp người quan trọng không kém hợp nghề của dòng họ.
\textit{(Career guidance should not rest only on family tradition. Suiting the person matters as much as continuing the family trade.)}
\end{baihoc}

\begin{ghinhoanh}
\textbf{Short English to remember:}

Career guidance should not rest only on family tradition.

Suiting the person matters as much as continuing the family trade.

\end{ghinhoanh}

\begin{tuhocnhanh}
\textbf{Cau hoi tu hoc / Self-study questions}

1. Van de chinh la gi? \textit{What is the main problem?}

2. Cach xu ly tot hon la gi? \textit{What is a better action?}

3. Mot cau tieng Anh em muon nho la cau nao? \textit{Which English sentence do you want to remember?}
\end{tuhocnhanh}
\ngancach

\section{Nghề Danh Giá Và Nghề Đứng Được}
\begin{truyen}{Nghề Danh Giá Và Nghề Đứng Được}{Prestigious Work and Work That Can Actually Sustain Life}
\chuhoa{C}{ó những nghề nghe sang tai hơn nghề khác.}
\textit{(Some professions sound more prestigious than others.)}

Nhưng nghề nghe sang chưa chắc là nghề em làm nổi và sống nổi.
\textit{(But a prestigious-sounding job is not necessarily one you can do well or live by.)}

Một học sinh mê ngành danh tiếng nhưng không chịu được kiểu công việc hàng ngày của nó.
\textit{(One student loved the status of a famous profession but could not endure its daily work style.)}

Ngược lại, một nghề ít hào nhoáng hơn lại vừa với tính khí và sức bền của em, nên em có thể đi rất xa.
\textit{(Meanwhile, a less glamorous field matched the student’s temperament and endurance, making a long path possible.)}

Danh tiếng là thứ xã hội nhìn từ ngoài.
\textit{(Prestige is what society sees from the outside.)}

Độ hợp là thứ em phải sống với nó từ trong.
\textit{(Fit is what you must live with from the inside.)}

\end{truyen}

\begin{baihoc}
Khi chọn nghề, đừng chỉ hỏi người khác sẽ nghĩ gì. Hãy hỏi mình có sống nổi với nhịp điệu thật của nghề đó không.
\textit{(When choosing work, do not ask only what others will think. Ask whether you can actually live with the true rhythm of that profession.)}
\end{baihoc}

\begin{ghinhoanh}
\textbf{Short English to remember:}

When choosing work, do not ask only what others will think.

Ask whether you can actually live with the true rhythm of that profession.

\end{ghinhoanh}

\begin{tuhocnhanh}
\textbf{Cau hoi tu hoc / Self-study questions}

1. Van de chinh la gi? \textit{What is the main problem?}

2. Cach xu ly tot hon la gi? \textit{What is a better action?}

3. Mot cau tieng Anh em muon nho la cau nao? \textit{Which English sentence do you want to remember?}
\end{tuhocnhanh}
\ngancach

\section{Người Học Việc Biết Mình Không Hợp}
\begin{truyen}{Người Học Việc Biết Mình Không Hợp}{The Apprentice Who Realized It Did Not Fit}
\chuhoa{M}{ột cậu học nghề theo lời khuyên của người lớn.}
\textit{(A boy apprenticed in a trade because adults advised it.)}

Sau nửa năm, em nhận ra mình không ghét vất vả, nhưng ghét đúng kiểu công việc đó.
\textit{(After half a year, he realized he did not hate hard work, but he did hate that specific kind of work.)}

Ban đầu em xấu hổ vì nghĩ bỏ là kém cỏi.
\textit{(At first he felt ashamed, thinking quitting meant weakness.)}

Sau này ông chủ cũ nói: ``Biết không hợp sớm còn hơn cố thêm mấy năm rồi hỏng luôn.''
\textit{(Later his former employer said, ``Knowing early that it does not fit is better than forcing yourself for years and breaking down.'')}

Rời một hướng không phải lúc nào cũng là thất bại.
\textit{(Leaving one direction is not always failure.)}

Có khi đó là một quyết định tiết kiệm cả tuổi trẻ.
\textit{(Sometimes it saves an entire youth.)}

Xã hội hay khen kiên trì, nhưng ít nói tới chuyện phải kiên trì đúng chỗ.
\textit{(Society praises perseverance, but says less about the need to persevere in the right place.)}

\end{truyen}

\begin{baihoc}
Không phải đường nào đã bắt đầu cũng phải đi tới cùng. Học sinh cần biết phân biệt giữa kiên định và cố chấp.
\textit{(Not every path once begun must be finished. Students need to distinguish steadfastness from stubbornness.)}
\end{baihoc}

\begin{ghinhoanh}
\textbf{Short English to remember:}

Not every path once begun must be finished.

Students need to distinguish steadfastness from stubbornness.

\end{ghinhoanh}

\begin{tuhocnhanh}
\textbf{Cau hoi tu hoc / Self-study questions}

1. Van de chinh la gi? \textit{What is the main problem?}

2. Cach xu ly tot hon la gi? \textit{What is a better action?}

3. Mot cau tieng Anh em muon nho la cau nao? \textit{Which English sentence do you want to remember?}
\end{tuhocnhanh}
\ngancach

\section{Một Nghề Cần Tay, Một Nghề Cần Mặt, Một Nghề Cần Tim}
\begin{truyen}{Một Nghề Cần Tay, Một Nghề Cần Mặt, Một Nghề Cần Tim}{One Job Needs Hands, One Needs Presence, One Needs Heart}
\chuhoa{C}{ó nghề chủ yếu dùng sức tay.}
\textit{(Some jobs rely mainly on manual skill.)}

Có nghề cần giao tiếp.
\textit{(Some require social presence.)}

Có nghề đòi hỏi ngồi rất lâu với chi tiết.
\textit{(Some demand long hours with detail.)}

Có nghề cần chịu áp lực liên tục.
\textit{(Some require constant pressure tolerance.)}

Một học sinh chỉ nhìn tên nghề mà không nhìn dạng lao động thật của nghề thì rất dễ chọn sai.
\textit{(A student who sees only the title of a job and not the true form of labor inside it is very likely to choose badly.)}

Cùng là giỏi, nhưng kiểu giỏi phù hợp với nghề nào mới là điều quyết định.
\textit{(Ability matters, but the type of ability that matches the job is what decides the matter.)}

Định hướng tốt phải cụ thể: em thích làm việc với người hay với vật?
\textit{(Good guidance must become concrete: do you like working with people or things?)}

thích chuyển động hay ngồi sâu?
\textit{(movement or deep stillness?)}

chịu được áp lực công khai hay chỉ hợp việc lặng lẽ?
\textit{(public pressure or quiet labor?)}

\end{truyen}

\begin{baihoc}
Tên nghề là lớp vỏ. Kiểu lao động mới là ruột. Học sinh cần được hướng vào phần ruột đó.
\textit{(The job title is the shell. The mode of labor is the substance. Students need guidance toward that substance.)}
\end{baihoc}

\begin{ghinhoanh}
\textbf{Short English to remember:}

The job title is the shell.

The mode of labor is the substance.

Students need guidance toward that substance.

\end{ghinhoanh}

\begin{tuhocnhanh}
\textbf{Cau hoi tu hoc / Self-study questions}

1. Van de chinh la gi? \textit{What is the main problem?}

2. Cach xu ly tot hon la gi? \textit{What is a better action?}

3. Mot cau tieng Anh em muon nho la cau nao? \textit{Which English sentence do you want to remember?}
\end{tuhocnhanh}
\ngancach

\section{Nghề Hot Hôm Nay, Nghề Nào Còn Lại Ngày Mai?}
\begin{truyen}{Nghề Hot Hôm Nay, Nghề Nào Còn Lại Ngày Mai?}{What Is Hot Today, and What Remains Tomorrow?}
\chuhoa{M}{ỗi thời có một danh sách nghề được tung hô.}
\textit{(Every era has a list of celebrated professions.)}

Nhưng thị trường thay nhanh hơn lời quảng cáo.
\textit{(But markets change faster than advertisements.)}

Một học sinh chọn ngành chỉ vì nghe nó đang hot.
\textit{(A student chose a major only because it was fashionable.)}

Ba năm sau, em phát hiện cạnh tranh đã khác và mình cũng không còn hứng thú như lúc đầu.
\textit{(Three years later, competition had changed and the student no longer felt the original excitement.)}

Đi theo xu hướng mà không có nền lực riêng giống như cưỡi sóng mà không biết bơi.
\textit{(Following trends without building your own strength is like riding a wave without knowing how to swim.)}

Xu hướng đáng tham khảo, nhưng không thể thay việc hiểu mình và luyện kỹ năng thật.
\textit{(Trends are worth consulting, but they cannot replace self-knowledge and the building of real skill.)}

\end{truyen}

\begin{baihoc}
Nghề hot chỉ cho em biết thị trường đang ồn về đâu. Nó không tự động cho em biết mình nên sống đời nào.
\textit{(A hot profession tells you where the market is noisy. It does not automatically tell you which life you should live.)}
\end{baihoc}

\begin{ghinhoanh}
\textbf{Short English to remember:}

A hot profession tells you where the market is noisy.

It does not automatically tell you which life you should live.

\end{ghinhoanh}

\begin{tuhocnhanh}
\textbf{Cau hoi tu hoc / Self-study questions}

1. Van de chinh la gi? \textit{What is the main problem?}

2. Cach xu ly tot hon la gi? \textit{What is a better action?}

3. Mot cau tieng Anh em muon nho la cau nao? \textit{Which English sentence do you want to remember?}
\end{tuhocnhanh}
\ngancach

\section{Tấm Bản Đồ Nghề Nghiệp Đầu Tiên}
\begin{truyen}{Tấm Bản Đồ Nghề Nghiệp Đầu Tiên}{The First Career Map}
\chuhoa{M}{ột cô giáo yêu cầu học sinh làm bản đồ nghề nghiệp cá nhân: điểm mạnh, việc mình chịu được lâu, việc mình ghét, kiểu môi trường mong muốn, mức sống mong muốn.}
\textit{(A teacher asked students to build personal career maps: strengths, work they could sustain, work they disliked, preferred environments, and desired living standards.)}

Nhiều em lần đầu thấy chọn nghề không phải chỉ là chọn một cái tên trên giấy đăng ký.
\textit{(For many, it was the first time realizing that choosing a career is not only choosing a title on an application form.)}

Nó là chọn một kiểu ngày tháng, một kiểu mệt, một kiểu niềm vui, một kiểu đánh đổi.
\textit{(It is choosing a pattern of days, a pattern of fatigue, a pattern of joy, and a pattern of trade-offs.)}

Người chọn nghề sớm chưa chắc giỏi hơn người chọn nghề chậm.
\textit{(Someone who chooses early is not necessarily wiser than someone who chooses slowly.)}

Nhưng người nghĩ nghiêm túc chắc chắn ít mù mờ hơn.
\textit{(But someone who thinks seriously is certainly less blind.)}

\end{truyen}

\begin{baihoc}
Học sinh cần công cụ để nghĩ về nghề cụ thể hơn, chứ không chỉ khẩu hiệu về đam mê và thành công.
\textit{(Students need tools to think about work concretely, not only slogans about passion and success.)}
\end{baihoc}

\begin{ghinhoanh}
\textbf{Short English to remember:}

Students need tools to think about work concretely, not only slogans about passion and success.

\end{ghinhoanh}

\begin{tuhocnhanh}
\textbf{Cau hoi tu hoc / Self-study questions}

1. Van de chinh la gi? \textit{What is the main problem?}

2. Cach xu ly tot hon la gi? \textit{What is a better action?}

3. Mot cau tieng Anh em muon nho la cau nao? \textit{Which English sentence do you want to remember?}
\end{tuhocnhanh}
